
\section{Method}

\subsection{Overview}

EAT replaces the standard MHSA block with a version that:

\begin{enumerate}
\item Uses \textbf{frozen base projections} transferred from the pretrained attention.
\item Attaches \textbf{trainable LoRA adapters} to each projection.
\item Optionally applies \textbf{progressive token pruning} at inference.

\end{enumerate}
\textbf{Figure 4.1} shows the pipeline.

\begin{lstlisting}
    +-------------------------------------------+
    |  Input X $\in$ R^{B$$\times$L$$\times$d}                       |
    +----------------+--------------------------+
                     |
                      
    +-------------------------------------------+
    |  Q/K/V projections (frozen W + LoRA)      |
    +----------------+--------------------------+
                     |
                      
    +-------------------------------------------+
    |  Optional RoPE                            |
    |  Optional GQA repeat_kv                   |
    +----------------+--------------------------+
                     |
                      
    +-------------------------------------------+
    |  Scaled dot-product attention             |
    |  + causal / padding mask                  |
    +----------------+--------------------------+
                     |
                      
    +-------------------------------------------+
    |  Output projection (frozen W + LoRA)      |
    +----------------+--------------------------+
                     |
                      
    +-------------------------------------------+
    |  Optional: Progressive token pruning       |
    |  (inference only, not training)           |
    +----------------+--------------------------+
                     |
                      
    Output Y
\end{lstlisting}

\subsection{Projections with LoRA}

\subsubsection{Base Weights}

The base projections $W_Q, W_K, W_V, W_O$ are \textbf{frozen}. They are 
transferred from the pretrained attention via \texttt{copy\_weights} ( 5).

\subsubsection{LoRA Adapters}

Each projection has trainable LoRA adapters:

\[
W_\text{eff} = W + \frac{\alpha}{r} B A
\tag{4.1}
\]

with $A \in \mathbb{R}^{r \times d_\text{in}}$, $B \in \mathbb{R}^{d_\text{out} \times r}$, 
$B = 0$ at initialization.

For the four projections:

\begin{table}[H]
\centering
\small
\begin{tabular}{llll}
\toprule
\textbf{Projection} & \textbf{Shape} & \textbf{Base frozen} & \textbf{LoRA added} \\
\midrule
$W_Q$ & $d_h H \times d$ & Yes & Yes \\
$W_K$ & $d_h H_{kv} \times d$ & Yes & Yes \\
$W_V$ & $d_h H_{kv} \times d$ & Yes & Yes \\
$W_O$ & $d \times d_h H$ & Yes & Yes \\
\bottomrule
\end{tabular}
\end{table}

\subsubsection{Why LoRA?}

LoRA provides three benefits:

\begin{enumerate}
\item \textbf{Parameter efficiency.} Only $r(d_\text{in} + d_\text{out})$ trainable
  parameters per projection.
\item \textbf{Identity at initialization.} Since $B = 0$, the initial forward pass
  matches the pretrained attention exactly.
\item \textbf{Recovery capacity.} LoRA can absorb distributional drift after
  injection.

\end{enumerate}
\subsection{Attention Computation}

The EAT forward pass is standard scaled dot-product attention:

\[
\text{head}_h = \text{softmax}\!\left(\frac{Q_h K_h^\top}{\sqrt{d_h}}\right) V_h
\tag{4.2}
\]

with $Q = W_Q X + \text{LoRA}_Q(X)$, etc.

\textbf{Causal mask.} For decoders:

\[
M_{ij} = \begin{cases} 0 & j \leq i \\ -\infty & j > i \end{cases}
\tag{4.3}
\]

\textbf{Padding mask.} For all architectures:

\[
M_{ij} = M_{ij} - \infty \cdot \mathbb{1}[\text{pad}(j)]
\tag{4.4}
\]

\subsection{Progressive Token Pruning}

\subsubsection{Motivation}

At inference time, the sequence length $L$ drives cost. If we can reduce $L$ 
for later layers, we save compute. The observation is that not all tokens 
are equally important for the final prediction.

\subsubsection{The Pruning Operator}

Given the output of an attention block $Y \in \mathbb{R}^{B \times L \times d}$ 
and a pruning module $\pi$:

\[
Y' = \pi(Y)
\tag{4.5}
\]

where $\pi$ reduces the sequence dimension from $L$ to $k \leq L$.

\subsubsection{Pruning Only at Inference}

Pruning is \textbf{disabled during training}:

\begin{lstlisting}
if self.prune_layer is not None and not self.training:
    out, _ = self.prune_layer(out)
\end{lstlisting}

This avoids gradient flow issues through \texttt{argsort}.

\subsubsection{The Pruning Module}

EAT delegates pruning to a \texttt{ProgressiveTokenPruning} module (see companion 
documentation). The module learns token importance scores and retains the 
top-$k$ tokens.

\subsection{Weight Transfer}

EAT relies on exact weight transfer from the pretrained attention. We 
support four cases:

\subsubsection{DistilBERT}

\begin{lstlisting}
target.q_lin.weight.copy_(src.q_lin.weight)
target.k_lin.weight.copy_(src.k_lin.weight)
target.v_lin.weight.copy_(src.v_lin.weight)
target.out_lin.weight.copy_(src.out_lin.weight)
\end{lstlisting}

\texttt{is\_nested\_encoder = False} -- DistilBERT applies \texttt{out\_lin} internally.

\subsubsection{RoBERTa}

\begin{lstlisting}
target.q_lin.weight.copy_(src.self.query.weight)
target.k_lin.weight.copy_(src.self.key.weight)
target.v_lin.weight.copy_(src.self.value.weight)
# out_proj skipped -- outer `attention.output.dense` applies it
\end{lstlisting}

\texttt{is\_nested\_encoder = True}.

\subsubsection{GPT-2}

\begin{lstlisting}
w = src.c_attn.weight.t()  # Conv1D transpose
q_w, k_w, v_w = w.chunk(3, dim=0)
target.q_lin.weight.copy_(q_w)
target.k_lin.weight.copy_(k_w)
target.v_lin.weight.copy_(v_w)
target.out_lin.weight.copy_(src.c_proj.weight.t())
\end{lstlisting}

\subsubsection{Qwen}

Direct copy from \texttt{q\_proj/k\_proj/v\_proj/o\_proj}. RoPE handled via \texttt{rotary\_emb} 
preservation.

\subsection{Silent Failure Modes}

EAT's \texttt{copy\_weights} is deliberately defensive: on failure, it logs a 
warning and \textbf{continues with random initialization}. This design was chosen 
to prevent injection from crashing, but it has a \textbf{critical implication}:

> If \texttt{copy\_weights} silently fails, the attention block keeps random 
> weights, and the model's forward pass produces garbage. Recovery 
> fine-tuning cannot recover from this because gradient signals are 
> dominated by noise.

We encountered this failure mode in practice (see  8.2). We discuss 
mitigation in  5.8.

\subsection{The Full Algorithm}

\textbf{Algorithm 1:} EAT forward pass
\begin{lstlisting}
Input:  X $\in$ R^{B$$\times$L$$\times$d}, attention_mask, position_ids
Output: Y $\in$ R^{B$$\times$L'$$\times$d}

1:  Q $\leftarrow$ W_Q X + LoRA_Q(X)           Frozen W + trainable LoRA
2:  K $\leftarrow$ W_K X + LoRA_K(X)
3:  V $\leftarrow$ W_V X + LoRA_V(X)
4:  if RoPE: (Q, K) $\leftarrow$ Rotary(Q, K)
5:  if GQA:  K, V $\leftarrow$ repeat_kv(K, V)
6:  S $\leftarrow$ Q K^T / $\sqrt{}$d_h
7:  S $\leftarrow$ S + causal_mask + padding_mask
8:  P $\leftarrow$ softmax(S)
9:  A $\leftarrow$ P V
10: A $\leftarrow$ concat_heads(A)
11: Y $\leftarrow$ W_O A + LoRA_O(A)
12: if inference and prune_layer:
13:     Y $\leftarrow$ prune_layer(Y)
14: return Y
\end{lstlisting}

\subsection{Complexity}

\begin{table}[H]
\centering
\small
\begin{tabular}{lll}
\toprule
\textbf{Component} & \textbf{Time} & \textbf{Space} \\
\midrule
Q/K/V projections & $\mathcal{O}(L d^2)$ & $\mathcal{O}(L d)$ \\
Attention & $\mathcal{O}(L^2 d)$ & $\mathcal{O}(L^2)$ \\
Output projection & $\mathcal{O}(L d^2)$ & $\mathcal{O}(L d)$ \\
LoRA (all) & $\mathcal{O}(L d r)$ & $\mathcal{O}(L r)$ \\
Pruning & $\mathcal{O}(L \log L)$ & $\mathcal{O}(L)$ \\
\textbf{Total (baseline)} & $\mathcal{O}(L^2 d)$ & $\mathcal{O}(L^2)$ \\
\textbf{Total (EAT)} & $\mathcal{O}(L^2 d + L d r)$ & $\mathcal{O}(L^2 + L r)$ \\
\bottomrule
\end{tabular}
\end{table}

\textbf{At inference with pruning}, $L$ is replaced by $k$ in downstream layers:

\begin{table}[H]
\centering
\small
\begin{tabular}{ll}
\toprule
\textbf{Component} & \textbf{Time} \\
\midrule
Downstream attention & $\mathcal{O}(k^2 d)$ \\
Downstream FFN & $\mathcal{O}(k d^2)$ \\
\bottomrule
\end{tabular}
\end{table}

For $k = L / 2$, downstream cost is reduced by 4$$\times$.

\subsection{Parameter Analysis}

EAT \textbf{does not reduce} parameters. For each projection of shape 
$d_\text{in} \times d_\text{out}$:

\begin{table}[H]
\centering
\small
\begin{tabular}{ll}
\toprule
\textbf{Configuration} & \textbf{Parameters} \\
\midrule
Baseline & $d_\text{in} d_\text{out}$ \\
EAT with LoRA rank $r$ & $d_\text{in} d_\text{out} + r(d_\text{in} + d_\text{out})$ \\
\bottomrule
\end{tabular}
\end{table}

For $d = 768$ and $r = 16$:

\begin{table}[H]
\centering
\small
\begin{tabular}{lll}
\toprule
\textbf{Config} & \textbf{Params} & \textbf{$\Delta$} \\
\midrule
Q projection baseline & 589,824 & -- \\
Q projection EAT & 589,824 + 24,576 = 614,400 & +24,576 \\
\bottomrule
\end{tabular}
\end{table}

\textbf{Across 4 projections $$\times$ $N_\ell$ layers:}

\begin{table}[H]
\centering
\small
\begin{tabular}{lll}
\toprule
\textbf{Model} & \textbf{Layers} & \textbf{EAT param increase} \\
\midrule
DistilBERT & 6 & +589,824 \\
RoBERTa & 12 & +1,179,648 \\
GPT-2 & 12 & +1,179,648 \\
\bottomrule
\end{tabular}
\end{table}

This matches the empirical observations in  7.
\begin{lstlisting}

---
\end{lstlisting}